\section{The Iterative Solver Based on ADMM Framework}
\label{section:admm}
In this section, we will present the methodology for solving the nonlinear system arising at each time step from the convex splitting scheme. 
We first detail the full derivation and algorithm of our proposed iterative solver based on ADMM framework. 
To benchmark its performance, we also implement the high-dimensional Newton's method for solving the same nonlinear system.

\subsection{First-order convex splitting numerical scheme}
For a given time $T>0$, consider a uniform partition of time with time step $\tau=\frac{T}{M}$ where $M\in \mathbb{N}_+$, that is, $0=t_0<t_1<\cdots<t_M=T$. 
The numerical solution at time $t_n=n\tau$ is denoted by $u^n$. 
\par The energy functional admits a decomposition into the sum of a convex and a concave part, i.e.
\begin{equation}
    E(u)=\int_{\Omega}(\frac{\epsilon^2}{2}|\nabla u|^2+u\log(u)+(1-u)\log(1-u))\text{d}x\text{d}y
    +\int_{\Omega}\theta (u-u^2)\text{d}x\text{d}y.
\end{equation}
\par The convex splitting scheme is constructed by treating the convex part implicitly and the concave part explicitly \cite{eyre1998unconditionally}, thus the scheme is stated as follows: 
\par For a given $u^n\in \ch$, find $u^{n+1}\in \ch$ such that
\begin{equation}
\label{scheme:cs}
    \frac{u^{n+1}-u^n}{\tau}-\epsilon^2\Delta_hu^{n+1}
    +\log(u^{n+1})-\log(1-u^{n+1})+\theta(1-2u^n)=0,
\end{equation}
which can also be written as\begin{equation}
\label{eq:cs}
    (\frac{1}{\tau}-\epsilon^2\Delta_h)u^{n+1}+\log(u^{n+1})-\log(1-u^{n+1})=\frac{1}{\tau}u^{n}-\theta(1-2u^n).
\end{equation}
The remaining problem is how to solve this nonlinear system efficiently. In the following, we present our algorithmic approach.
\subsection{Iterative solver based on ADMM framework}
Motivated by the need for robust solvers, our main strategy is to reinterpret the numerical scheme \eqref{scheme:cs} 
through an optimization problem according to the theory of convex optimization in \cite{boyd2004convex}, which enables us to leverage powerful tools from convex optimization.
\par For a given $u^n\in \ch$, we define a discrete energy functional
\begin{equation}
    J(u)=\frac{1}{2\tau}\|u-u^n\|_h^2+\frac{\epsilon^2}{2}\|\nabla_hu\|_h^2+
    \langle u,\log(u)\rangle_h+
    \langle 1-u,\log(1-u)\rangle_h
    +\theta \langle u,1-2u^n \rangle_h,
\end{equation}
over the set $u\in \ch$ \cite{xu2019stability}. 
Then the numerical scheme \eqref{scheme:cs} is mathematically equivalent to the following optimization problem,
\begin{equation}\label{eq:opt}
    u^{n+1}=\arg \min_{u\in \ch}J(u).
\end{equation}
 We then decompose the energy functional $J(u)$ into a sum of two convex terms, that is,
 \begin{subequations}\label{eq:Ju}
     \begin{align}
         &J(u)=J_1(u)+J_2(u),\\ 
         &J_1(u)=\frac{1}{2\tau}\|u-u^n\|_h^2+\frac{\epsilon^2}{2}\|\nabla_hu\|_h^2,\\
         &J_2(u)=\langle u,\log(u)\rangle_h+
    \langle 1-u,\log(1-u)\rangle_h
    +\theta \langle u,1-2u^n \rangle_h,
     \end{align}
 \end{subequations}
and now formalize the optimization problem by putting it into the following standard formulation, which reveals its structure clearly.
\begin{mini!}
{u_1, u_2}{J_1(u_1)+J_2(u_2)}{}{}
\addConstraint{u_1}{=u_2}{}
\addConstraint{u_i}{\in \ch,}{\quad i=1,2}
\end{mini!}
\par Accordingly, the optimization problem is posed in terms of its Lagrange function, defined by
\begin{equation}\label{eq:Lagrange}
    \mathcal{L}(u_1,u_2;u_3)=J_1(u_1)+J_2(u_2)+
    \langle u_3,u_1-u_2 \rangle_h,
\end{equation}
where $u_3 \in \ch$ is the Lagrange multiplier \cite{rockafellar1973multiplier}. And the corresponding augmented Lagrange function is
\begin{equation}
    \mathcal{L}_{\rho}(u_1,u_2;u_3)=\mathcal{L}(u_1,u_2;u_3)+\frac{\rho}{2}\|u_1-u_2\|_h^2,
\end{equation}
where $\rho>0$ is the penalty parameter. We employ the standard ADMM framework \cite{neal2011distributed} to the optimization problem, whose $k$-th iteration step is given by
\begin{subequations}
    \begin{align}
        \label{eq:subp1}
        &u_2^{(k+1)}=\arg \min_{u_2\in \ch} \mathcal{L}_{\rho}(u_1^{(k)},u_2;u_3^{(k)}),\\
        \label{eq:subp2}
        &u_1^{(k+1)}=\arg \min_{u_1\in \ch} \mathcal{L}_{\rho}(u_1,u_2^{(k+1)};u_3^{(k)}),\\
        &u_3^{(k+1)}=u_3^{(k)}+\alpha\rho(u_1^{(k+1)}-u_2^{(k+1)}),
    \end{align}
\end{subequations}
where $\alpha\in(0,\frac{\sqrt{5}+1}{2})$ is the step size for multiplier update.
\par The efficiency of our ADMM implementation hinges on solving the two subproblems in \eqref{eq:subp1} and \eqref{eq:subp2}. 
We solve them by exploiting their distinct mathematical structures.
\par The first-order optimality condition for the $u_2$-update is given  by
\begin{equation}
    \label{eq:u2}
    \log(u_2^{(k+1)})-\log(1-u_2^{(k+1)})+
    \theta(1-2u^n)-u_3^{(k)}-\rho(u_1^{(k)}-u_2^{(k+1)})=0,
\end{equation}                                                            
Specifically, the optimality condition \eqref{eq:u2} reduces to solving a set of independent, one-dimensional nonlinear equations at each grid point $\{(i,j)|0\leq i,j \leq N-1\}$, 
completely decoupling the global problem into $N^2$ independent scalar equations. 
This allows us to solve all $N^2$ equations in parallel, each using a few steps of a safeguarded one-dimensional Newton's method \cite{dennis1996numerical,ortega2000iterative}.

\par The $u_1$-update results in a symmetric positive-definite linear system, i.e., 
    \begin{equation}
    \label{eq:u1}
        (\frac{1}{\tau}-\epsilon^2\Delta_h+\rho)u_1^{(k+1)}=\frac{1}{\tau}u^n+\rho u_2^{(k+1)}-u_3^{(k)},
    \end{equation}
which we solve with a preconditioned conjugate gradient (PCG) method \cite{concus1976generalized}. 
For uniform grids, this system is circulant and can be solved via discrete fast Fourier transform (FFT) techniques \cite{cooley1965algorithm}.

\par To summarize, we present the detailed process of the iterative solver based on ADMM framework in \Cref{alg:ADMM}.
\begin{algorithm}[!t]
    \caption{Iterative solver based on ADMM framework}
\label{alg:ADMM}
    \raggedright
    1. Define $u_1^{(0)}=u^n$, $u_2^{(0)}=u^n$, and
    $u_3^{(0)}=0$.\\
    2. Select $\rho\in(0,+\infty), \alpha \in (0,\frac{\sqrt{5}+1}{2})$.\\
    3. For $k\in\mathbb{N}$, update $u_2$ by solving the nonlinear system
    \begin{equation}
    \log(u_2^{(k+1)})-\log(1-u_2^{(k+1)})+
    \theta(1-2u^n)-u_3^{(k)}-\rho(u_1^{(k)}-u_2^{(k+1)})=0,
    \end{equation}
    \hspace{10pt} for $u_2^{(k+1)}\in\ch$ with one-dimensional Newton's method at each grid point.\\
    4. For $k\in\mathbb{N}$, update $u_1$ by solving the linear system
    \begin{equation}
    \label{u1}
        (\frac{1}{\tau}-\epsilon^2\Delta_h+\rho)u_1^{(k+1)}=\frac{1}{\tau}u^n+\rho u_2^{(k+1)}-u_3^{(k)},
    \end{equation}
    \hspace{10pt} for $u_1^{(k+1)}\in\ch$ with PCG or FFT.\\ 
    5. Update the Lagrange multipliers $u_3^{(k+1)}\in\ch$ as
    \begin{equation}\label{eq:u3}
    u_3^{(k+1)}=u_3^{(k)}+\alpha \rho(u_1^{(k+1)}-u_2^{(k+1)}).
    \end{equation}
\end{algorithm}
\subsection{Iterative solver based on high-dimensional Newton's method}
In contrast to the optimization-based approach, we may alternatively view the nonlinear system \eqref{eq:cs} directly as a root-finding problem 
and tackle it with a high-dimensional Newton iteration \cite{ortega2000iterative}. 
This leads directly to the canonical Newton iteration for systems of equations, whose $k$-th iteration step is given by
\begin{equation}\label{eq:newton1}
    \begin{aligned}
    &(\frac{1}{\tau}-\epsilon^2\Delta_h+G^{(k)})(u^{(k+1)}-u^{(k)})\\
    =&-(\frac{1}{\tau}-\epsilon^2\Delta_h)u^{(k)}-\log(u^{(k)})
    +\log(1-u^{(k)})
    +\frac{1}{\tau}u^n-\theta(1-2u^n),
\end{aligned}
\end{equation}
where $G^{(k)}$ is the Jacobian matrix of the nonlinear function $g(u^{(k)})=\log(u^{(k)})-\log(1-u^{(k)})$ for $u^{(k)}\in \ch$. 
It follows directly from its definition that $G^{(k)}$ is a diagonal matrix with strictly positive diagonal entries. 
In fact, $G^{(k)}_{j,j}=\frac{1}{u_{j}^{(k)}(1-u_{j}^{(k)})}(0\leq j \leq N^2-1)$ when we treat $u^{(k)} \in \ch$ as a $N^2$-dimensional vector.
\par The $k$-iteration step \eqref{eq:newton1} also results in a symmetric positive-definite linear system, i.e.,
\begin{equation}
\begin{aligned}
    &(\frac{1}{\tau}-\epsilon^2\Delta_h+G^{(k)})u^{(k+1)}\\
    =&G^{(k)}u^{(k)}
    -\log(u^{(k)})
    +\log(1-u^{(k)})
    +\frac{1}{\tau}u^n-\theta(1-2u^n),
\end{aligned}
\end{equation}
which we solve with PCG method \cite{concus1976generalized}. Due to the lack of circulant structure in the coefficient matrix, the discrete FFT is not directly applicable. 
\par Similarly, we present the process of the iterative solver based on high-dimensional Newton's method in \Cref{alg:Newton}.

\begin{algorithm}[!t]
    \caption{Iterative solver based on high-dimensional Newton's method}
\label{alg:Newton}
    \raggedright
    1. Define $u^{(0)}=u^n$.\\
    2. For $k\in\mathbb{N}$, update $u$ by solving the linear system
    \begin{equation}
\begin{aligned}
    &(\frac{1}{\tau}-\epsilon^2\Delta_h+G^{(k)})u^{(k+1)}\\
    =&G^{(k)}u^{(k)}
    -\log(u^{(k)})
    +\log(1-u^{(k)})
    +\frac{1}{\tau}u^n-\theta(1-2u^n),
\end{aligned}
\end{equation}
    \hspace{10pt} for $u^{(k+1)}\in\ch$ with PCG.\\
    
\end{algorithm}
    